\subsubsection{多维势的向量化扩展}\label{app:vector-potential}
若引入向量势
\[
\boldsymbol{\varphi}(e)=\bigl(\varphi_1(e),\dots,\varphi_d(e)\bigr),
\qquad
B(\boldsymbol{\theta})_{ij}=\sum_{e:i\to j}\exp\!\bigl(\langle\boldsymbol{\theta},\boldsymbol{\varphi}(e)\rangle\bigr),
\]
则压力
\[
P(\boldsymbol{\theta})=\log\rho\bigl(B(\boldsymbol{\theta})\bigr)
\]
与大偏差率函数
\[
I(\alpha)=\sup_{\boldsymbol{\theta}\in\RR^d}\bigl(\langle\boldsymbol{\theta},\alpha\rangle-(P(\boldsymbol{\theta})-P(0))\bigr)
\]
均保持"有限矩阵 + 行列式"的可审计口径；二维/三维实例见后续小节。

\subsubsection{负携带势的显式闭式（可审计实例）}
取
\[
Q_-:=\{0\overline{1}2,\,1\overline{1}2,\,01\overline{1},\,11\overline{1}\},
\qquad
\phi(s\to t):=\mathbf{1}_{\{t\in Q_-\}}.
\]
令 $u=e^\theta$。此时 $B_\theta$ 等价于把落入 $Q_-$ 的列整体乘上 $u$，并有显式因式分解
\[
\det(I-zB_\theta)
\;=\;
\bigl(u z^3-u z^2+u z+1\bigr)\bigl(3u z^3+u z^2-u z-2z^2-2z+1\bigr).
\]
主极点 $z_\star(\theta)=\lambda(\theta)^{-1}$ 由第二因子给出；代入 $z=1/\lambda$ 得三次方程
\[
\lambda^3-(u+2)\lambda^2+(u-2)\lambda+3u=0,\qquad u=e^\theta,
\]
从而压力 $P(\theta)=\log\lambda(\theta)$ 完全由该三次方程的最大实根确定。

\leanverified{paper\_vector\_potential\_phi\_minus\_ldp\_param}
\begin{proposition}[负携带势的代数大偏差：斜率的有理参数化与闭式速率函数]\label{prop:vector-potential-phi-minus-ldp-param}
沿用本小节的记号。令 \(u=e^{\theta}>0\)，并令 \(\lambda(u)\) 为三次方程
\[
F(\lambda,u):=\lambda^3-(u+2)\lambda^2+(u-2)\lambda+3u=0
\]
在 \(u>0\) 上的 Perron 分支（最大正实根）。
定义压力 \(P(\theta)=\log\lambda(e^\theta)\) 与斜率 \(\alpha(\theta):=P'(\theta)\)，并以 Legendre 变换定义率函数
\[
I(\alpha):=\sup_{\theta\in\RR}\bigl(\theta\alpha-(P(\theta)-P(0))\bigr).
\]
则存在参数 \(\lambda\in(1+\sqrt3,\infty)\) 的显式有理参数化，使得
\[
u(\lambda)=\frac{\lambda(\lambda^{2}-2\lambda-2)}{\lambda^{2}-\lambda-3},
\qquad
\alpha(\lambda)=\frac{\lambda^{4}-3\lambda^{3}-3\lambda^{2}+8\lambda+6}{\lambda^{4}-2\lambda^{3}-5\lambda^{2}+12\lambda+6},
\]
并且在该参数表示下率函数具有闭式
\[
I(\alpha(\lambda))=\alpha(\lambda)\log u(\lambda)-\log\lambda+\log 3.
\]
\end{proposition}
\begin{proof}
由 \(F(\lambda,u)=0\) 视为关于 \(u\) 的线性方程，直接解得
\[
u=\frac{\lambda^{3}-2\lambda^{2}-2\lambda}{\lambda^{2}-\lambda-3}
=\frac{\lambda(\lambda^{2}-2\lambda-2)}{\lambda^{2}-\lambda-3}
=:u(\lambda).
\]
对 \(u=e^{\theta}\) 由隐函数微分得
\[
\lambda_u=-\frac{\partial_u F}{\partial_\lambda F},
\qquad
\alpha(\theta)=\frac{\dd}{\dd\theta}\log\lambda(e^\theta)=\frac{u\,\lambda_u}{\lambda}
=-\frac{u\,\partial_u F(\lambda,u)}{\lambda\,\partial_\lambda F(\lambda,u)}\Bigg|_{\lambda=\lambda(u)}.
\]
对本例，
\(
\partial_u F(\lambda,u)=-(\lambda^2-\lambda-3)
\)
与
\(
\partial_\lambda F(\lambda,u)=3\lambda^2-2(u+2)\lambda+(u-2)
\)。
代入并用 \(u=u(\lambda)\) 化简，即得到所示有理式 \(\alpha(\lambda)\)。
\par
对率函数，在极值点满足 \(\alpha=P'(\theta)\) 的口径下有
\[
I(\alpha)=\theta\alpha-P(\theta)+P(0).
\]
此处 \(\theta=\log u(\lambda)\)、\(P(\theta)=\log\lambda\)。并且当 \(u=1\) 时
\(
F(\lambda,1)=\lambda^{3}-3\lambda^{2}-\lambda+3=(\lambda-3)(\lambda^{2}-1)
\)，故 Perron 根为 \(\lambda(1)=3\)，从而 \(P(0)=\log 3\)。
代入即得
\(
I(\alpha(\lambda))=\alpha(\lambda)\log u(\lambda)-\log\lambda+\log 3
\)。
证毕。
\end{proof}

由隐式微分或脚本数值（\texttt{--phi neg\_state}）可得
\[
P'(0)=\frac{1}{8}\ \ (\approx 0.1250000006),\qquad
P''(0)=\frac{25}{128}\ \ (\approx 0.1953128548).
\]
进一步，由三次方程在 \(\theta=0\) 处的解析分支可把更高阶累积量（偏斜、峰度等）闭式导出；下式由脚本 \texttt{scripts/exp\_sync\_kernel\_phi\_minus\_cubic\_series.py} 一键生成并可回归测试：

% Generated by scripts/exp_sync_kernel_phi_minus_cubic_series.py
\[
\begin{aligned}
P(\theta)
&=\frac{2471 \theta^{5}}{524288} + \frac{19705 \theta^{4}}{786432} + \frac{59 \theta^{3}}{1024} + \frac{25 \theta^{2}}{256} + \frac{\theta}{8} + \log{\left(3 \right)}+O(\theta^{6}),\\
P''(0)&=\frac{25}{128},\qquad P^{(3)}(0)=\frac{177}{512},\qquad P^{(4)}(0)=\frac{19705}{32768},\qquad P^{(5)}(0)=\frac{37065}{65536},\\
\gamma_1&:=\frac{P^{(3)}(0)}{(P''(0))^{3/2}}=\frac{354 \sqrt{2}}{125}\ \ (\approx 4.00505),\\
\gamma_2&:=\frac{P^{(4)}(0)}{(P''(0))^2}=\frac{3941}{250}\ \ (\approx 15.764).
\end{aligned}
\]

因此典型频率 $\alpha_\star=P'(0)=1/8$，局部二次近似为
\[
I(\alpha)\approx \frac{(\alpha-\frac18)^2}{2P''(0)}=\frac{64}{25}\Bigl(\alpha-\frac18\Bigr)^2,
\]
并且端点代价满足
\[
I(0)=\log\frac{3}{1+\sqrt3},\qquad I(1)=\log 3.
\]
加权 primitive 轨道谱 $p_n(\theta)$ 仍由 Möbius 反演给出，并满足
\[
p_n(\theta)\asymp \frac{e^{nP(\theta)}}{n},
\qquad
M_{B,\theta}(N):=\sum_{n\le N}\frac{p_n(\theta)}{e^{nP(\theta)}}=\log N+\mathsf{M}_{B,\theta}+o(1),
\]
其中常数项 $\mathsf{M}_{B,\theta}$ 可由 Abel 有限部（$z\to 1/\lambda(\theta)$）读出。

\subsubsection{二维加权素数谱（输出位 + 负携带）}
定义二维势
\[
\varphi_e(e)=e,\qquad \varphi_-(s\to t)=\mathbf{1}_{\{t\in Q_-\}},
\]
并令
\[
u:=e^{\theta_e},\qquad v:=e^{\theta_-}.
\]
加权邻接矩阵条目为
\[
B(u,v)_{ij}=\sum_{e:i\to j}u^{\varphi_e(e)}v^{\varphi_-(e)}.
\]
对应行列式可化为 6 次多项式
\[
\begin{aligned}
\Delta(z;u,v)=&\,1-(u+1)z-u(v^2+2v+2)z^2\\
&+uv(uv+2u+v+2)z^3+uv^2(-u^2+3u-1)z^4\\
&+\bigl(u^4v^2-2u^3v^2-u^3v-2u^2v^2-u^2v+uv^2\bigr)z^5\\
&+u^2v^2(u^2+u+1)z^6,
\end{aligned}
\]
因此
\[
\zeta(z;u,v)=\frac{1}{\Delta(z;u,v)}.
\]
记 $z_\star(u,v)$ 为最小正根，$\lambda(u,v)=1/z_\star(u,v)$，压力
\[
P(\theta_e,\theta_-)=\log\lambda(e^{\theta_e},e^{\theta_-}).
\]
在 $(\theta_e,\theta_-)=(0,0)$ 处有
\[
\partial_{\theta_e}P(0,0)=\frac12\ \ (\approx 0.5000000000),\qquad
\partial_{\theta_-}P(0,0)=\frac18\ \ (\approx 0.1250000006),
\]
\[
\partial^2_{\theta_e\theta_e}P(0,0)=\frac{11}{102}\ \ (\approx 0.1078436229),\quad
\partial^2_{\theta_-\theta_-}P(0,0)=\frac{25}{128}\ \ (\approx 0.1953128548),\quad
\partial^2_{\theta_e\theta_-}P(0,0)=0.
\]
因此局部二次近似为
\[
P(\theta_e,\theta_-)=\log 3+\frac12\theta_e+\frac18\theta_-+\frac{11}{204}\theta_e^2+\frac{25}{256}\theta_-^2+O(|\theta|^3).
\]
二维大偏差率函数
\[
I(\alpha_e,\alpha_-)=\sup_{\theta_e,\theta_-}\bigl(\theta_e\alpha_e+\theta_-\alpha_- - (P(\theta_e,\theta_-)-P(0,0))\bigr)
\]
在典型点附近满足
\[
I(\alpha_e,\alpha_-)\approx \frac{51}{11}\Bigl(\alpha_e-\frac12\Bigr)^2+\frac{64}{25}\Bigl(\alpha_- -\frac18\Bigr)^2.
\]
加权 primitive 轨道数
\[
p_n(\theta_e,\theta_-)\asymp \frac{e^{nP(\theta_e,\theta_-)}}{n}
\]
并由 Möbius 反演从 $a_n(\theta_e,\theta_-)=\mathrm{Tr}(B(u,v)^n)$ 提取。数值复现见脚本 \texttt{scripts/exp\_sync\_kernel\_weighted\_pressure\_2d.py}。

\begin{proposition}[二维势的二阶严格解耦与首个非高斯混合项]\label{prop:sync-kernel-2d-exact-decouple}
\leanverified{paper\_sync\_kernel\_2d\_exact\_decouple}
对本小节的二维势 \((\varphi_e,\varphi_-)\)，压力
\(
P(\theta_e,\theta_-)=\log\lambda(e^{\theta_e},e^{\theta_-})
\)
满足以下精确结构结论：
\begin{enumerate}
  \item \textbf{（二阶严格解耦）}
  \[
  \boxed{\ \partial^2_{\theta_e\theta_-}P(0,0)=0\ }.
  \]
  因而在典型点附近的二次型近似中\textbf{无交叉项}，二维局部率函数二次项严格可分离。
  \item \textbf{（三阶“单向耦合”）}
  \[
  \boxed{\ \partial^3_{\theta_e\theta_-\theta_-}P(0,0)=0,\qquad
  \partial^3_{\theta_e\theta_e\theta_-}P(0,0)=-\frac{1453}{36992}<0\ }.
  \]
  \item \textbf{（四阶混合曲率为负）}
  \[
  \boxed{\ \partial^4_{\theta_e\theta_e\theta_-\theta_-}P(0,0)=-\frac{422483}{10061824}<0,\qquad
  \partial^4_{\theta_e\theta_-\theta_-\theta_-}P(0,0)=0\ }.
  \]
\end{enumerate}
\end{proposition}

\begin{proof}
这里不需要数值差分。因为 \(\Delta(z;u,v)\) 是显式整系数多项式且
\(
z_\star(1,1)=1/3
\)
对应 \(\lambda(1,1)=3\)，可在 \((u,v)=(1,1)\) 邻域把最小正根写为解析分支
\(
z_\star(u,v)
\)
并用隐式方程 \(\Delta(z_\star(u,v);u,v)=0\) 对 \((\theta_e,\theta_-)\) 作隐式微分（或等价地作幂级数消元），得到所示导数值。
脚本 \texttt{scripts/exp\_sync\_kernel\_weighted\_pressure\_2d\_exact\_mixed\_derivatives.py} 对该推导做了完全有理数的自动化核验并输出可直接 \verb|\input| 的 \LaTeX\ 片段：
% Auto-generated; do not edit by hand.
\begin{equation}\label{eq:sync-kernel-weighted-pressure-2d-exact-mixed-derivatives}
\boxed{
\begin{aligned}
\partial_{\theta_e\theta_-}^2 P(0,0) &= 0,\\
\partial_{\theta_e\theta_e\theta_-}^3 P(0,0) &= -\frac{1453}{36992},\qquad\partial_{\theta_e\theta_-\theta_-}^3 P(0,0) = 0,\\
\partial_{\theta_e\theta_e\theta_-\theta_-}^4 P(0,0) &= -\frac{422483}{10061824},\qquad\partial_{\theta_e\theta_-\theta_-\theta_-}^4 P(0,0) = 0.
\end{aligned}
}
\end{equation}

证毕。
\end{proof}

\begin{remark}[局部率函数的“严格可分离二次项”与可检验的非高斯耦合指纹]\label{rem:sync-kernel-2d-nongaussian-coupling}
命题 \ref{prop:sync-kernel-2d-exact-decouple} 的第一条给出
\(
\nabla^2 P(0,0)
\)
在 \((\theta_e,\theta_-)\) 坐标下的严格对角结构，因此
\[
I(\alpha_e,\alpha_-)
=\frac{(\alpha_e-\frac12)^2}{2\,\partial^2_{\theta_e\theta_e}P(0,0)}
\;+\;
\frac{(\alpha_--\frac18)^2}{2\,\partial^2_{\theta_-\theta_-}P(0,0)}
\;+\;O(|\alpha-\alpha_\star|^3),
\]
且二次项\emph{严格无交叉项}。
另一方面，三、四阶混合导数给出首个非高斯层面的耦合常数：例如三阶项
\(
\partial^3_{\theta_e\theta_e\theta_-}P(0,0)<0
\)
预测了在长样本下
\(
(\bar\varphi_e-\tfrac12)^2(\bar\varphi_- -\tfrac18)
\)
的中心化平均出现负号偏斜；四阶混合曲率为负则给出“联合尾部”相对于独立二次模型的偏离方向。
\end{remark}

\subsubsection{三维加权素数谱（输出位 + 负携带 + 碰撞）}
定义三维势
\[
\varphi_e(e)=e,\qquad
\varphi_-(s\to t)=\mathbf{1}_{\{t\in Q_-\}},\qquad
\varphi_2(d)=\mathbf{1}_{\{d=2\}},
\]
并令
\[
u:=e^{\theta_e},\qquad v:=e^{\theta_-},\qquad w:=e^{\theta_2}.
\]
按状态顺序 $Q$，三维加权邻接矩阵为
\[
B(u,v,w)=
\begin{pmatrix}
1 & 1 & w & 0 & 0 & 0 & 0 & 0 & 0 & 0\\
0 & 0 & 0 & 1 & 1 & w & 0 & 0 & 0 & 0\\
u & uw & 0 & 0 & 0 & 0 & 0 & 0 & 0 & v\\
0 & 0 & 0 & 0 & 1 & 1 & uvw & 0 & 0 & 0\\
u & u & uw & 0 & 0 & 0 & 0 & 0 & 0 & 0\\
0 & 0 & 0 & u & u & uw & 0 & 0 & 0 & 0\\
0 & 0 & 0 & 1 & w & 0 & 0 & 0 & v & 0\\
0 & 0 & 0 & u & uw & 0 & 0 & 0 & uv & 0\\
0 & 1 & 1 & 0 & 0 & 0 & 0 & vw & 0 & 0\\
0 & u & u & 0 & 0 & 0 & 0 & uvw & 0 & 0
\end{pmatrix}.
\]
对应行列式仍为 6 次多项式
\[
\Delta(z;u,v,w)=1+c_1z+c_2z^2+c_3z^3+c_4z^4+c_5z^5+c_6z^6,
\]
其中
\[
c_1=-(uw+1),\qquad
c_2=-u(v^2w+vw+v+2),\qquad
c_3=uv(uw+1)(vw+w+1),
\]
\[
c_4=-u\bigl(u^2v^2w^3-3uv^2w+uv(w^2+w-2)+u(w^2-w)+v^2w\bigr),
\]
\[
c_5=uvw\bigl(u^3 v w^3-2u^2 v w-u^2 w^3-u v(w+1)-u+v\bigr),
\]
\[
c_6=u^2 v^2 w(2w-1)(u^2w+u+1).
\]
因此三观测量的加权 $\zeta$-函数具有有理闭式
\[
\zeta(z;u,v,w)=\frac{1}{\Delta(z;u,v,w)}.
\]
\leanverified{paper\_sync\_kernel\_3d\_critical\_w\_half}
\begin{proposition}[内部临界超平面 \texorpdfstring{$w=\frac12$}{w=1/2}：有效谱维数的一阶坍缩]\label{prop:sync-kernel-3d-critical-w-half}
对任意 \(u>0,v>0\)，当 \(w=\tfrac12\)（等价于 \(\theta_2=\log w=-\log 2\)）时，有
\[
c_6(u,v,\tfrac12)=0,
\]
从而
\(
\Delta(z;u,v,\tfrac12)
\)
从 \(6\) 次降为至多 \(5\) 次多项式。
在本文“非零谱维数为 \(6\)”的接口口径下，这等价于：\(B(u,v,\tfrac12)\) 至少额外产生一个结构性零特征值，
其非零谱维数从 \(6\) 降到至多 \(5\)。
\end{proposition}
\begin{proof}
由显式公式 \(c_6=u^2v^2w(2w-1)(u^2w+u+1)\) 直接代入 \(w=\tfrac12\) 得 \(c_6=0\)。证毕。
\end{proof}

\leanverified{paper\_sync\_kernel\_3d\_recurrence\_drop}
\begin{corollary}[周期迹递推阶数的骤降（可审计指纹）]\label{cor:sync-kernel-3d-recurrence-drop}
令 \(a_n(u,v,w):=\Tr(B(u,v,w)^n)\)。
在 \(w=\tfrac12\) 处，序列 \(a_n(u,v,\tfrac12)\) 存在阶数至多 \(5\) 的线性递推关系；
而在一般正权重域 \(w\neq\tfrac12\) 时该接口给出的递推阶数为 \(6\)。
\end{corollary}

\leanverified{paper\_sync\_kernel\_3d\_critical\_quintic}
\begin{proposition}[临界五次式指纹：\texorpdfstring{$(u,v,w)=(1,1,\tfrac12)$}{(u,v,w)=(1,1,1/2)}]\label{prop:sync-kernel-3d-critical-quintic}
在 \((u,v,w)=(1,1,\tfrac12)\) 处，\(\Delta(z;1,1,\tfrac12)\) 的 \(z^6\) 项严格消失，并具有可复现的五次式输出：
% Auto-generated by scripts/exp_sync_kernel_weighted_3d_exact_audits.py
\[
\boxed{\;
\Delta\!\left(z;1,1,\tfrac12\right)
=1-\frac{3}{2}z-4z^2+3z^3+\frac{19}{8}z^4-\frac{5}{4}z^5,
\;}
\qquad(c_6=0).
\]

\end{proposition}
记 $\lambda(u,v,w)=1/z_\star(u,v,w)$，压力
\[
P(\theta_e,\theta_-,\theta_2)=\log\lambda(e^{\theta_e},e^{\theta_-},e^{\theta_2}).
\]
在 $\theta=0$ 处有
\[
\nabla P(0)=\left(\frac12,\frac18,\frac13\right)
\ \ (\approx (0.5000000000,\ 0.1250000006,\ 0.3333333335)),
\]
并且 Hessian 满足
\[
\nabla^2P(0)=
\begin{pmatrix}
\frac{11}{102} & 0 & \frac{181}{1632}\\
0 & \frac{25}{128} & \frac{5}{128}\\
\frac{181}{1632} & \frac{5}{128} & \frac{2}{9}
\end{pmatrix},
\]
即 $\partial^2_{\theta_e\theta_-}P(0)=0$，而 $\partial^2_{\theta_e\theta_2}P(0),\partial^2_{\theta_-\theta_2}P(0)>0$。因此
\[
P(\theta)=\log 3+\frac12\theta_e+\frac18\theta_-+\frac13\theta_2
+\frac{11}{204}\theta_e^2+\frac{25}{256}\theta_-^2+\frac19\theta_2^2
+\frac{181}{1632}\theta_e\theta_2+\frac{5}{128}\theta_-\theta_2+O(|\theta|^3).
\]
这意味着在二阶近似下 \(\varphi_e\) 与 \(\varphi_-\) \textbf{零相关}，但二者分别与 \(\varphi_2\) 正相关；因此在以 \(\varphi_2\) 为中介作条件化后会出现\textbf{条件负相关}（Schur 补读出）：

% Generated by scripts/exp_sync_kernel_3d_conditional_covariance.py
\[
\begin{aligned}
\mathrm{Cov}(\varphi_e,\varphi_-\mid\varphi_2)=\partial_{\theta_e\theta_-}^2P(0)-\frac{\partial_{\theta_e\theta_2}^2P(0)\,\partial_{\theta_-\theta_2}^2P(0)}{\partial_{\theta_2\theta_2}^2P(0)}
&=- \frac{2715}{139264}\ \ (\approx -0.0194953),\\
\mathrm{Corr}(\varphi_e,\varphi_-\mid\varphi_2)=\frac{\mathrm{Cov}(\varphi_e,\varphi_-\mid\varphi_2)}{\sqrt{\mathrm{Var}(\varphi_e\mid\varphi_2)\,\mathrm{Var}(\varphi_-\mid\varphi_2)}}
&\approx -0.196016.
\end{aligned}
\]


\noindent
从而“零相关但非独立”在本三势耦合中具有一个可审计的定量证书。
三维大偏差率函数
\[
I(\alpha)=\sup_{\theta\in\mathbb{R}^3}\bigl(\theta\cdot\alpha-(P(\theta)-P(0))\bigr)
\]
在典型点 $\alpha_\star=(\frac12,\frac18,\frac13)$ 附近满足
\[
I(\alpha)\approx \tfrac12(\alpha-\alpha_\star)^\top\bigl(\nabla^2P(0)\bigr)^{-1}(\alpha-\alpha_\star),
\]
其中
\[
\bigl(\nabla^2P(0)\bigr)^{-1}\approx
\begin{pmatrix}
19.8119 & 2.04960 & -10.2480\\
2.04960 & 5.51860 & -1.99298\\
-10.2480 & -1.99298 & 9.96491
\end{pmatrix}.
\]
\begin{proposition}[三势 Hessian 的精确逆矩阵（有理数证书）]\label{prop:sync-kernel-3d-hessian-inverse-exact}
\leanverified{paper\_sync\_kernel\_3d\_hessian\_inverse\_exact}
上式的逆矩阵可写成完全精确的有理数矩阵：
% Auto-generated by scripts/exp_sync_kernel_weighted_3d_exact_audits.py
\[
\boxed{\;
\bigl(\nabla^2P(0)\bigr)^{-1}=
\begin{pmatrix}
\frac{1713192}{86473} & \frac{886176}{432365} & -\frac{886176}{86473}\\
\frac{886176}{432365} & \frac{2386048}{432365} & -\frac{861696}{432365}\\
-\frac{886176}{86473} & -\frac{861696}{432365} & \frac{861696}{86473}
\end{pmatrix}
\;}
\]

\end{proposition}

\begin{corollary}[二阶解耦基：\texorpdfstring{$(2\!+\!1)$}{2+1} 正交分量的显式构造]\label{cor:sync-kernel-3d-orthobasis}
\leanverified{paper\_sync\_kernel\_3d\_orthobasis}
令 \(\varphi_e,\varphi_-,\varphi_2\) 为本小节三条边势。由 \(\partial^2_{\theta_e\theta_-}P(0)=0\) 与 Hessian 的显式值，
可构造一个可审计的二阶正交基 \((\eta,\psi',\varphi_2)\)：
% Auto-generated by scripts/exp_sync_kernel_weighted_3d_exact_audits.py
\[
\begin{aligned}
\eta&:=\varphi_e-\frac{543}{1088}\,\varphi_2,\\
\psi&:=\varphi_- -\frac{45}{256}\,\varphi_2,\\
\psi'&:=\psi+\frac{27693}{74564}\,\eta,\\
\mathrm{Var}(\eta)&=\frac{93205}{1775616},\\
\mathrm{Var}(\psi')&=\frac{432365}{2386048},\\
\mathrm{Var}(\varphi_2)&=\frac{2}{9}.
\end{aligned}
\]

其在 \(\theta=0\) 的二次近似下两两不相关（协方差为零），从而二次型可被严格对角化。
\end{corollary}

\begin{remark}[三通道势—响应辛相空间与 \texorpdfstring{$\mathrm{Sp}(6)$}{Sp(6)} 各向同性扫描（可复现）]\label{rem:sync-kernel-3d-sp6-isotropy-scan}
在三维加权势的口径下，压力 \(P(\theta_e,\theta_-,\theta_2)\) 给出自然的“势—响应”共轭对：
\[
J_e:=\partial_{\theta_e}P,\qquad
J_-:=\partial_{\theta_-}P,\qquad
J_2:=\partial_{\theta_2}P.
\]
因此相空间坐标 \((\theta;J)\in\RR^6\) 携带标准辛形式
\[
\omega=\dd\theta_e\wedge \dd J_e+\dd\theta_-\wedge \dd J_-+\dd\theta_2\wedge \dd J_2,
\]
其保持群为 \(\mathrm{Sp}(6)\)（见注 \ref{rem:sp6-minimal-continuous-envelope} 的统一解释口径）。
在二次近似层，涨落由 Fisher/协方差矩阵
\(
\Sigma(\theta):=\nabla^2P(\theta)
\)
控制；在最小块对角提升 \(C=\mathrm{diag}(\Sigma,\Sigma)\) 下，\(C\) 的辛特征值退化为 \(\Sigma\) 的三个特征值，
从而“单尺度（各向同性）”可用一个完全可审计的标量分数度量：
\[
\mathcal{I}(\theta):=\log\frac{\lambda_{\max}(\Sigma(\theta))}{\lambda_{\min}(\Sigma(\theta))}.
\]
脚本 \texttt{\detokenize{scripts/exp_sync_kernel_weighted_sp6_isotropy_scan.py}} 对 \((u,v,w)=(e^{\theta_e},e^{\theta_-},e^{\theta_2})\) 作网格扫描并输出最优候选点（最小化 \(\mathcal{I}\)），其生成表格见表 \ref{tab:sync_kernel_weighted_sp6_isotropy_scan}。
\end{remark}
\begin{table}[H]
\centering
\scriptsize
\setlength{\tabcolsep}{6pt}
\caption{Sp(6)-isotropy scan for the 3D weighted sync-kernel pressure. At each grid point $(u,v,w)$ we compute $\Sigma=\nabla^2P(\theta)$ (Fisher/covariance in $\theta=(\log u,\log v,\log w)$) by central finite differences, and report the best candidates minimizing $\log(\lambda_{\max}/\lambda_{\min})$ over eigenvalues of $\Sigma$.}
\label{tab:sync_kernel_weighted_sp6_isotropy_scan}
\begin{tabular}{r r r r r}
\toprule
$u$ & $v$ & $w$ & eigs$(\Sigma)$ & $\log(\lambda_{\max}/\lambda_{\min})$\\
\midrule
0.8 & 0.8 & 0.3 & (0.03781,0.06239,0.1634) & 1.46388\\
0.7 & 0.8 & 0.3 & (0.03691,0.06057,0.1599) & 1.46606\\
0.9 & 0.8 & 0.3 & (0.03828,0.06407,0.1668) & 1.47179\\
0.7 & 0.9 & 0.3 & (0.03836,0.06827,0.1681) & 1.47752\\
0.8 & 0.9 & 0.3 & (0.03898,0.07088,0.1716) & 1.48212\\
0.6 & 0.8 & 0.3 & (0.03543,0.05852,0.1562) & 1.48327\\
0.6 & 0.9 & 0.3 & (0.03718,0.06541,0.1642) & 1.4854\\
1 & 0.8 & 0.3 & (0.03845,0.06566,0.17) & 1.4864\\
0.9 & 0.9 & 0.3 & (0.03922,0.07331,0.1748) & 1.4946\\
0.9 & 0.7 & 0.3 & (0.03613,0.05603,0.1612) & 1.49549\\
0.8 & 0.7 & 0.3 & (0.03531,0.05518,0.1577) & 1.49631\\
0.6 & 0.8 & 0.4 & (0.04132,0.06593,0.1846) & 1.49696\\
0.5 & 0.8 & 0.4 & (0.04007,0.06209,0.1791) & 1.49767\\
1 & 0.7 & 0.3 & (0.03662,0.05684,0.1646) & 1.5031\\
1.1 & 0.8 & 0.3 & (0.03842,0.06715,0.1731) & 1.50534\\
0.6 & 0.7 & 0.4 & (0.03945,0.05739,0.1782) & 1.5076\\
0.7 & 0.7 & 0.3 & (0.03407,0.0542,0.154) & 1.50882\\
0.7 & 0.7 & 0.4 & (0.04043,0.05949,0.1834) & 1.51186\\
1 & 0.9 & 0.3 & (0.03921,0.07558,0.1778) & 1.51198\\
0.7 & 0.8 & 0.4 & (0.04176,0.06934,0.1894) & 1.51208\\
\bottomrule
\end{tabular}
\end{table}


由此得到一个自然的 arity-charge
\[
\chi:=\varphi_e-\varphi_2\in\{-1,0,1\},
\]
其典型密度与方差密度为
\[
\mathbb{E}[\chi]=\frac12-\frac13=\frac16,\qquad
\mathrm{Var}_\infty(\chi)=\frac{11}{102}+\frac{2}{9}-2\cdot\frac{181}{1632}=\frac{265}{2448}\ (\approx 0.1082516340).
\]
三维加权 primitive 轨道数
\[
p_n(\theta_e,\theta_-,\theta_2)=\sum_{\substack{\tau\ \mathrm{primitive}\|\tau|=n}}
\exp\!\bigl(\theta_e S_\tau\varphi_e+\theta_-S_\tau\varphi_-+\theta_2S_\tau\varphi_2\bigr)
\]
仍满足 $p_n(\theta)\asymp e^{nP(\theta)}/n$，并可由
$a_n(\theta)=\mathrm{Tr}(B(u,v,w)^n)$ 经 Möbius 反演提取。数值复现见脚本 \texttt{scripts/exp\_sync\_kernel\_weighted\_pressure\_3d.py}。
在未加权情形 $(u,v,w)=(1,1,1)$ 下，有
\[
(a_1,\dots,a_{10})=(2,14,20,94,222,770,2116,6694,19442,59494),
\]
\[
(p_1,\dots,p_{10})=(2,6,6,20,44,123,302,825,2158,5926),
\]
并呈现典型 prime-orbit 主项 $p_n\sim 3^n/n$。

\subsubsection{多维 Big-Witt 同余护栏的零成本推广：系数环升级不改变整除/同余结构}\label{app:multi-dim-witt-guardrails}
本附录的二维/三维势并不改变 primitive--periodic 的 Big-Witt/necklace 结构：其差异仅在系数环从 \(\ZZ[u]\) 升级到多变量多项式环。
因此，附录 \ref{app:witt-bridge-package} 的整除律与 Dwork 型同余可逐字推广为多维版本，并立即给出一个无需额外计算开销且可直接复算的“实现自审计护栏”。

\paragraph{二维势：\texorpdfstring{$R=\ZZ[u,v]$}{R=Z[u,v]}}
对二维加权矩阵 \(B(u,v)\)（\S\ref{app:vector-potential}），令
\[
a_n(u,v):=\Tr\bigl(B(u,v)^n\bigr)\in\ZZ[u,v],
\qquad
\psi^d:\ZZ[u,v]\to\ZZ[u,v],\ \ \psi^d(f)(u,v):=f(u^d,v^d).
\]
则由命题 \ref{prop:witt-necklace-dwork-R}（取 \(R=\ZZ[u,v]\)）直接推出：
\begin{enumerate}
  \item \textbf{（多维 necklace 整除）}对任意 \(n\ge 1\)，
  \[
  \boxed{\ \sum_{d\mid n}\mu(d)\,a_{n/d}(u^d,v^d)\in n\,\ZZ[u,v]\ }.
  \]
  \item \textbf{（多维 Dwork 同余）}对任意素数 \(p\) 与 \(k\ge 1\)，
  \[
  \boxed{\ a_{p^k}(u,v)\equiv a_{p^{k-1}}(u^p,v^p)\pmod{p^k}\quad\text{于}\ \ZZ[u,v]\ }.
  \]
\end{enumerate}

\paragraph{三维势：\texorpdfstring{$R=\ZZ[u,v,w]$}{R=Z[u,v,w]}}
对三维加权矩阵 \(B(u,v,w)\)（见本附录三维势小节给出的显式 \(10\times 10\) 矩阵），令
\[
a_n(u,v,w):=\Tr\bigl(B(u,v,w)^n\bigr)\in\ZZ[u,v,w],
\qquad
\psi^d(f)(u,v,w):=f(u^d,v^d,w^d).
\]
同理得到：
\begin{enumerate}
  \item \textbf{（多维 necklace 整除）}对任意 \(n\ge 1\)，
  \[
  \boxed{\ \sum_{d\mid n}\mu(d)\,a_{n/d}(u^d,v^d,w^d)\in n\,\ZZ[u,v,w]\ }.
  \]
  \item \textbf{（多维 Dwork 同余）}对任意素数 \(p\) 与 \(k\ge 1\)，
  \[
  \boxed{\ a_{p^k}(u,v,w)\equiv a_{p^{k-1}}(u^p,v^p,w^p)\pmod{p^k}\quad\text{于}\ \ZZ[u,v,w]\ }.
  \]
\end{enumerate}

\begin{remark}[自审计要点：只用迹序列做模同余检查]\label{rem:multi-dim-self-audit}
上述两条约束对任何“声称来自 Euler 乘积/primitive 分解”的多维核都是必要条件。
因此无需谱分解、无需先算 primitive：只要能计算 \(a_n(\cdot)=\Tr(B(\cdot)^n)\)，就可在系数环内做
\(\bmod n\) 的整除审计与 \(\bmod p^k\) 的 Frobenius 兼容审计；一旦失败，即可判定实现不可能来自本文的 \(\zeta\)--Euler 乘积体系。
\end{remark}
